%% ============================================================
%% 动态时序评价段落模板
%% 变量:
%%   n_periods            — 时段数量
%%   period_names         — 时段名称列表 (如 ["2019","2020","2021"])
%%   object_names         — 评价对象名称列表
%%   n_samples            — 对象数量
%%   eval_method          — 单期评价方法 (如 "TOPSIS")
%%   time_weights         — 时间权重列表
%%   time_weight_method   — 时间权重确定方法说明
%%   period_scores        — 分期得分矩阵 [period][object]
%%   dynamic_scores       — 综合动态得分列表
%%   dynamic_ranks        — 综合排名列表
%%   trend_plot_path      — 趋势折线图路径
%%   dynamic_bar_path     — 动态综合得分柱状图路径
%% ============================================================

\subsection{动态时序综合评价模型}

\subsubsection{模型原理}

当评价数据具有时间序列特征（如多年度面板数据）时，需构建动态评价模型以捕捉各对象的发展趋势。

本文采用"分期评价 + 时间加权汇总"的策略：

\textbf{步骤一：分期评价。} 对每个时段 $t \in \{<< period_names | join(",\\ ") >>\}$，独立采用 << eval_method >> 方法计算各对象的得分 $s_i(t)$。

\textbf{步骤二：确定时间权重。} << time_weight_method >>，得到时间权重向量 $\boldsymbol{\tau} = (<< time_weights | map("%.4f" | format) | join(",\\ ") >>)$。通常赋予较近时段更高权重以反映"近期优先"原则。

\textbf{步骤三：动态综合得分。}
\begin{equation}
    S_i^{dyn} = \sum_{t=1}^{T} \tau_t \cdot s_i(t)
    \label{eq:dynamic_score}
\end{equation}

\subsubsection{分期评价结果}

\begin{table}[H]
    \centering
    \caption{各时段 << eval_method >> 得分}
    \label{tab:dynamic_period_scores}
    \begin{tabular}{l<% for p in period_names %>c<% endfor %>}
        \toprule
        \textbf{评价对象} & <% for p in period_names %><% if not loop.first %> & <% endif %>\textbf{<< p >>}<% endfor %> \\
        \midrule
<% for i in range(n_samples) %>
        << object_names[i] >> & <% for t in range(n_periods) %><% if not loop.first %> & <% endif %><< "%.4f" | format(period_scores[t][i]) >><% endfor %> \\
<% endfor %>
        \bottomrule
    \end{tabular}
\end{table}

\subsubsection{动态综合评价结果}

\begin{table}[H]
    \centering
    \caption{动态综合评价得分与排名}
    \label{tab:dynamic_result}
    \begin{tabular}{lcc}
        \toprule
        \textbf{评价对象} & \textbf{动态综合得分} & \textbf{排名} \\
        \midrule
<% for i in range(n_samples) %>
        << object_names[i] >> & << "%.6f" | format(dynamic_scores[i]) >> & << dynamic_ranks[i] >> \\
<% endfor %>
        \bottomrule
    \end{tabular}
\end{table}

<% if trend_plot_path %>
\begin{figure}[H]
    \centering
    \includegraphics[width=0.85\textwidth]{<< trend_plot_path >>}
    \caption{各评价对象得分趋势变化}
    \label{fig:dynamic_trend}
\end{figure}
<% endif %>

<% if dynamic_bar_path %>
\begin{figure}[H]
    \centering
    \includegraphics[width=0.75\textwidth]{<< dynamic_bar_path >>}
    \caption{动态综合得分排序}
    \label{fig:dynamic_bar}
\end{figure}
<% endif %>